When AI agents share channels and credentials, one agent can press another to break its operator's rule, a first step towards collusion.
Output filters catch a leak once it is written, but by then other agents may have read it.
Can an overseer tell beforehand that an agent is about to give in?
We show that a linear probe on the agent's activations can, in a controlled multi-agent setting.
Across six scenarios, a peer presses a target for up to 20 rounds to reveal a secret, such as a deploy token; the target often refuses, then gives in.
Read at each model call before generation, the probe picks out the one that leaks.
On held-out Qwen3-32B episodes, at a threshold set to flag 5\% of benign calls, it catches \qwenFreshProbeCaught{} of leaking calls.
At the same flag rate, a zero-shot decision-model monitor catches \qwenMatchedJevCaught{}.
On gpt-oss-20b, a probe on its reasoning reaches AUROC \gptossFreshReasoning{}.
A bag-of-words classifier trained on the same calls comes close, so most of the gain comes from training on labels, not activations.
In a counterfactual replay, allowing sharing leaves the probe's score unchanged: it reads the coming disclosure, not the violation.
The warning is short-range, and halting on any alarm would stop most benign episodes, which a second text check does not fix.
For multi-agent oversight, this suggests monitors that act before an agent gives in, where checking its output comes too late: messages others read immediately, paraphrased or encoded leaks, and recruitment spreading through a swarm.
